% TOP_PROBLEM: {"id":30007667,"problem_number":"LOCAL-30007667","title":"Elite capture under inclusive institutions","category_id":21,"category":{"id":21,"name":"economics","display_name":"Economics","description":"Open economic theory statements, published research agendas, and proposed scoped specifications; question provenance and historical priority are qualified per record.","slug":"economics","order_index":21,"created_at":"2026-10-08T00:00:00Z"},"status":"research_question","external_url":null,"legacy_ids":[],"published":false,"statement":"In the explicitly specified setting: How do economic elites preserve political influence after formal democratization, and which institutional changes prevent capture without simply creating new entrenched elites? The mathematical statement above fixes the target, assumptions and scope of this catalogue formulation.","economics":{"original_id":156,"field":"Political economy and institutions","kind":"identification","formalization_basis":"adapted_research_question","source_status":"research_agenda","priority_status":"earliest_found","priority_note":"Earliest directly inspected qualifying passage in this audit; earlier origins were not established. A review or research-agenda statement does not imply original priority.","earliest_verified_reference":{"authors":["World Bank"],"title":"World Development Report 2017: Governance and the Law","year":2017,"url":"https://www.worldbank.org/en/publication/wdr2017","doi":null,"locator":"About, opening three paragraphs; Main Messages, first three bullets","evidence_summary":"The report asks why beneficial policies fail adoption or implementation, emphasizing commitment, power asymmetries, capture, and coalition formation."},"current_reference":{"authors":["World Bank"],"title":"World Development Report 2017: Governance and the Law","year":2017,"url":"https://www.worldbank.org/en/publication/wdr2017","doi":null,"locator":"About, opening three paragraphs; Main Messages, first three bullets","evidence_summary":"The report asks why beneficial policies fail adoption or implementation, emphasizing commitment, power asymmetries, capture, and coalition formation."},"formalization_note":"The cited passage establishes a research direction, not this exact mathematical statement. The notation, population, policy menu, assumptions, and estimands are an original adaptation for this catalogue; no universal unsolved theorem or source priority is claimed."},"statusQualification":"Published question or research agenda verified at the cited locator. The mathematical specification is adapted for this catalogue and includes additional assumptions; source publication does not establish first-ever priority or certify this exact model remains unsolved. The cited passage establishes a research direction, not this exact mathematical statement. The notation, population, policy menu, assumptions, and estimands are an original adaptation for this catalogue; no universal unsolved theorem or source priority is claimed.","statusReviewedAt":"2026-10-08"}
% FIRST_POSTED: 2026-10-08
% ENTRY_KIND: statement_only
% !TeX program = lualatex
\documentclass[11pt]{article}
\usepackage{fontspec}
\usepackage[a4paper,margin=25mm]{geometry}
\usepackage{amsmath,amssymb,mathtools}
\usepackage{xurl}
\usepackage[hidelinks]{hyperref}
\newcommand{\sref}[2]{\hyperlink{source:#1:#2}{[S#2]}}
\setlength{\emergencystretch}{3em}
\begin{document}
% problemId:30007667
\section*{Elite capture under inclusive institutions}\label{30007667}
\subsection{Definitions and mathematical statement}
Fix a polity undergoing a precisely specified institutional reform \(z\), such as expanded electoral participation or competitive entry into political office. Let \(P_{it}(z)\) be actor \(i\)'s political influence, measured through offices, appointments, policy access, or networks according to a declared rule. Let \(E\) denote actors classified as incumbent elites using prereform information, and define elite influence share \(S_t(z)=\sum_{i\in E}P_{it}(z)/\sum_iP_{it}(z)\), with a positive denominator. Let \(Y_t(z)\) measure broad economic welfare. Estimate changes in \(S\), policy responsiveness, and \(Y\) over a horizon allowing elite adaptation. Decompose descriptive persistence into routes such as party control, financing, social networks, economic ownership, or relatives taking office, while causal channel claims require additional variation. Evaluate complementary reforms \(p\) that enhance contestability and prevent excluded citizens or new elites from recreating capture. Inclusive legal forms do not ensure inclusive policy outcomes. The task adapts the governance report's explicit problem of exclusion and capture; the specific influence metric, dynamic intervention, and replacement-elite test are our specification. Existing country studies offer partial answers, and original priority for this broad formulation is not established.

\par\textbf{Formulation and scope.} The cited passage establishes a research direction, not this exact mathematical statement. The notation, population, policy menu, assumptions, and estimands are an original adaptation for this catalogue; no universal unsolved theorem or source priority is claimed.

\par\textbf{Open-question provenance.} An explicit question or research agenda was verified at the source locator. This does not by itself certify that every later version remains unresolved \sref{30007667}{1}.

\par\textbf{Historical priority.} Earliest directly inspected qualifying passage in this audit; earlier origins were not established. A review or research-agenda statement does not imply original priority.

\subsection{Short English statement}
In the explicitly specified setting: How do economic elites preserve political influence after formal democratization, and which institutional changes prevent capture without simply creating new entrenched elites? The mathematical statement above fixes the target, assumptions and scope of this catalogue formulation.

\subsection{Sources}
\begin{itemize}\raggedright
\item[\textnormal{[S1]}]\hypertarget{source:30007667:1}{} World Bank. 2017. World Development Report 2017: Governance and the Law. About, opening three paragraphs; Main Messages, first three bullets.
\url{https://www.worldbank.org/en/publication/wdr2017}.
{\small\textit{Earliest verified question/agenda reference; historical priority qualified below. The report asks why beneficial policies fail adoption or implementation, emphasizing commitment, power asymmetries, capture, and coalition formation.}}
\end{itemize}
\end{document}
